\chapter{Bảo mật trong hệ thống mạng CAN}

\section{Các vấn đề về bảo mật trong hệ thống mạng CAN}
Mạng CAN (\textit{Controller Area Network}) được thiết kế từ đầu nhằm tối ưu cho tốc độ truyền thông và độ tin cậy trong các hệ thống nhúng, đặc biệt là ô tô. Tuy nhiên, thiết kế ban đầu của CAN không tập trung vào các yêu cầu bảo mật, vì vậy giao thức này tồn tại nhiều lỗ hổng tiềm ẩn khi được đặt trong môi trường kết nối hiện đại.

Dưới đây là một số vấn đề bảo mật chính của mạng CAN:

\begin{itemize}[leftmargin=*]
    \item \textbf{Không có cơ chế xác thực (Authentication):} Bất kỳ thiết bị nào được kết nối vào bus CAN đều có thể gửi thông điệp. Hệ thống không có cách tự nhiên để xác minh nguồn gốc thực sự của thông điệp đó.
    \item \textbf{Không có mã hóa (Encryption):} Dữ liệu trên bus CAN được truyền ở dạng rõ ràng (\textit{plaintext}), tạo điều kiện cho kẻ tấn công nghe lén và phân tích nội dung.
    \item \textbf{Không có phân quyền truy cập (Access Control):} Tất cả các node trên mạng đều chia sẻ cùng một môi trường truyền thông, không có cơ chế kiểm soát vai trò hay đặc quyền truy cập.
    \item \textbf{Dễ bị tấn công từ chối dịch vụ (Denial-of-Service - DoS):} Một thiết bị độc hại có thể liên tục gửi các frame có ID ưu tiên cao để chiếm quyền truy cập bus, làm gián đoạn hoạt động của các node khác.
    \item \textbf{Khả năng nghe trộm (Eavesdropping):} Vì toàn bộ bus chia sẻ chung tín hiệu, mọi node đều có thể quan sát các thông điệp đang lưu thông trên mạng.
    \item \textbf{Dễ bị giả mạo (Spoofing):} Một thiết bị độc hại có thể giả mạo ID của một node hợp lệ để phát thông điệp sai lệch lên mạng.
\end{itemize}

\section{Các phương pháp bảo mật trong hệ thống mạng CAN}
Để khắc phục các hạn chế nêu trên, nhiều giải pháp bảo mật đã được đề xuất. Mỗi phương pháp có ưu điểm và đánh đổi riêng, tùy thuộc vào khả năng phần cứng, yêu cầu thời gian thực và kiến trúc tổng thể của hệ thống.

\subsection{Xác thực thông điệp (Message Authentication)}
\textbf{Nguyên lý.} Một mã xác thực thông điệp (MAC - \textit{Message Authentication Code}) được gắn thêm vào dữ liệu để kiểm tra nguồn gốc và tính toàn vẹn của thông điệp.

\textbf{Cách thực hiện.}
\begin{itemize}[leftmargin=*]
    \item Sử dụng hàm băm kết hợp với khóa bí mật để tạo mã xác thực, ví dụ như HMAC-SHA256.
    \item Node gửi tính giá trị MAC từ dữ liệu và truyền kèm trong payload.
    \item Node nhận tính lại MAC và so sánh với giá trị đã nhận; nếu trùng khớp, thông điệp được xem là hợp lệ.
\end{itemize}

\textbf{Ví dụ.} Với dữ liệu $D$ và khóa $K$, node gửi tính $\text{MAC} = \text{HMAC}_{K}(D)$ rồi gửi cặp $(D, \text{MAC})$. Node nhận tính lại $\text{HMAC}_{K}(D)$ để kiểm tra.

\textbf{Ưu điểm.} Ngăn chặn giả mạo thông điệp. Phát hiện được các sửa đổi trái phép trong quá trình truyền.

\textbf{Nhược điểm.} Không tự chống được tấn công phát lại nếu không kết hợp thêm bộ đếm hoặc dấu thời gian. Làm tăng kích thước payload, dễ chạm giới hạn 8 byte của CAN 2.0. Tăng tải tính toán trên vi điều khiển.

\subsection{Sử dụng bộ đếm tin nhắn đơn giản (Simple Message Counter)}
\textbf{Nguyên lý.} Mỗi thông điệp được gắn một bộ đếm tăng dần để đảm bảo tính ``tươi mới'' (\textit{freshness}). Thiết bị nhận dùng giá trị này để phát hiện và từ chối các bản tin bị phát lại.

\textbf{Cách thực hiện.}
\begin{itemize}[leftmargin=*]
    \item Gắn một trường bộ đếm nhỏ, ví dụ 8 bit hoặc 16 bit, vào thông điệp $M_k$.
    \item Thiết bị nhận lưu giá trị bộ đếm gần nhất cho từng ID hoặc từng phiên giao tiếp.
    \item Nếu bộ đếm trong thông điệp mới không lớn hơn giá trị đã lưu, thông điệp có thể bị xem là bản sao hoặc lỗi thời.
\end{itemize}

\textbf{Ví dụ.} Nếu thông điệp trước đó có bộ đếm là 5 thì thông điệp kế tiếp hợp lệ phải có bộ đếm bằng 6 hoặc lớn hơn.

\textbf{Ưu điểm.} Cơ chế đơn giản, hiệu quả và phù hợp với vi điều khiển giới hạn tài nguyên. Giúp chống lại các cuộc tấn công phát lại ở mức cơ bản.

\textbf{Nhược điểm.} Không tự bảo vệ trước giả mạo hoặc nghe lén nếu không kết hợp với MAC hay mã hóa. Yêu cầu thiết bị nhận phải lưu trạng thái bộ đếm, có thể làm tăng nhu cầu bộ nhớ.

\subsection{Phát hiện xâm nhập (Intrusion Detection System - IDS)}
\textbf{Nguyên lý.} IDS giám sát và phân tích luồng dữ liệu trên bus CAN để phát hiện các hành vi bất thường hoặc độc hại.

\textbf{Cách thực hiện.}
\begin{itemize}[leftmargin=*]
    \item Thu thập dữ liệu CAN và đánh giá bất thường dựa trên tập luật (\textit{rule-based}) hoặc mô hình học máy.
    \item So sánh với hành vi bình thường để phát hiện dấu hiệu lạ như tần suất gửi tăng đột biến hoặc xuất hiện ID không mong đợi.
    \item Gửi cảnh báo hoặc kích hoạt cơ chế chặn khi phát hiện rủi ro.
\end{itemize}

\textbf{Ví dụ.} Nếu một ID thường chỉ gửi 10 thông điệp mỗi phút nhưng đột nhiên phát 100 thông điệp trong vài giây, IDS có thể xem đây là dấu hiệu tấn công.

\textbf{Ưu điểm.} Không đòi hỏi thay đổi sâu vào giao thức CAN gốc. Có khả năng phát hiện các kiểu tấn công mới chưa được định nghĩa trước.

\textbf{Nhược điểm.} Có thể tạo ra cảnh báo giả (\textit{false positive}). Yêu cầu năng lực xử lý tương đối cao nếu dùng mô hình phức tạp.

